% Part: first-order-logic
% Chapter: completeness
% Section: outline

\documentclass[../../../include/open-logic-section]{subfiles}

\begin{document}

\iftag{FOL}
      {\olfileid[id]{fol}{com}{out}}
      {\olfileid[id]{pl}{com}{out}}

\olsection{Garis Besar Pembuktian}

Pembuktian teorema kelengkapan agak rumit, dan pembaca mudah kehilangan
arah ketika pertama kali membacanya. Karena itu, mari kita uraikan garis
besar pembuktiannya. Langkah pertama ialah mengubah sudut pandang agar
kita dapat melihat jalan menuju suatu bukti. Jika kelengkapan dipahami
sebagai ``setiap kali $\Gamma \Entails !A$, berlaku $\Gamma \Proves
!A$,'' bahkan menemukan gagasannya saja mungkin sulit: untuk menunjukkan
bahwa $\Gamma \Proves !A$, kita harus menemukan !!a{derivation}, sedangkan
hipotesis $\Gamma \Entails !A$ tampaknya sama sekali tidak membantu kita
menemukannya. Pada beberapa sistem bukti, !!a{derivation} dapat dibangun
secara langsung, tetapi kita akan mengambil pendekatan yang sedikit
berbeda. Perubahan sudut pandang yang diperlukan adalah sebagai berikut:
kelengkapan juga dapat dirumuskan dengan pernyataan ``jika $\Gamma$
konsisten, maka himpunan itu dapat dipenuhi.'' Barangkali kita dapat
menggunakan informasi dalam~$\Gamma$, bersama dengan hipotesis bahwa
himpunan itu konsisten, untuk membangun
\iftag{FOL}{!!a{structure}}{!!a{valuation}} yang memenuhi setiap
\iftag{FOL}{!!{sentence}}{!!{formula}} dalam~$\Gamma$. Bagaimanapun, kita
mengetahui jenis \iftag{FOL}{!!{structure}}{!!{valuation}} yang kita
cari: yang persis seperti digambarkan oleh~$\Gamma$!{}

Jika $\Gamma$ hanya memuat \iftag{FOL}{!!{sentence}s
  atomik}{!!{propositional variable}s}, model untuknya mudah dibangun.
\iftag{FOL}{ Andaikan semua !!{sentence}s atomik itu berbentuk
  $\Atom{P}{a_1,\dots,a_n}$, dengan $a_i$ sebagai !!{constant}s.}{}
Kita hanya perlu menentukan
\iftag{FOL}{%
  !!a{domain}~$\Domain{M}$ dan suatu penetapan bagi~$P$ sedemikian
  sehingga $\Sat{M}{\Atom{P}{a_1,\ldots,a_n}}$. Namun, hal itu tidak
  terlalu sulit: tetapkan $\Domain{M} = \Nat$, $\Assign{\Obj c_i}{M} =
  i$, dan untuk setiap $\Atom{P}{a_1,\ldots,a_n} \in \Gamma$, masukkan
  tupel $\tuple{k_1, \dots, k_n}$ ke dalam $\Assign{P}{M}$, dengan $k_i$
  sebagai indeks simbol konstanta~$a_i$ (yakni, $a_i \ident \Obj
  c_{k_i}$).}{%
  !!a{valuation}~$\pAssign{v}$ sedemikian sehingga $\pSat{v}{p}$ untuk
  semua $p \in \Gamma$. Tetapkan saja $\pAssign{v}(p) = \True$ jika dan
  hanya jika $p \in \Gamma$.}

Sekarang andaikan $\Gamma$ memuat suatu !!{formula}~$\lnot !B$, dengan
$!B$ atomik. Kita mungkin khawatir bahwa konstruksi
$\iftag{FOL}{\Struct{M}}{\pAssign{v}}$ mengganggu kemungkinan untuk
membuat $\lnot !B$ benar. Namun, di sinilah konsistensi~$\Gamma$
berperan: jika $\lnot !B \in \Gamma$, maka $!B \notin \Gamma$; jika
tidak demikian, $\Gamma$ akan tak konsisten. Jika $!B \notin \Gamma$,
maka menurut konstruksi $\iftag{FOL}{\Struct{M}}{\pAssign{v}}$ kita,
$\iftag{FOL}{\Sat/{M}}{\pSat/{v}}{!B}$, sehingga
$\iftag{FOL}{\Sat{M}}{\pSat{v}}{\lnot !B}$. Sejauh ini tidak ada
masalah.

Bagaimana jika $\Gamma$ memuat formula kompleks yang tidak atomik?
Misalkan himpunan itu memuat $!A \land !B$. Agar formula itu benar, kita
harus bertindak seolah-olah $!A$ dan $!B$ keduanya berada
dalam~$\Gamma$. Jika $!A \lor !B \in \Gamma$, kita harus membuat
sekurang-kurangnya salah satunya benar, yakni bertindak seolah-olah salah
satunya berada dalam~$\Gamma$.

Hal ini menyarankan gagasan berikut: kita tambahkan !!{formula}s lain
ke~$\Gamma$ sedemikian rupa sehingga (a)~himpunan yang dihasilkan tetap
konsisten, (b)~untuk setiap !!{sentence} atomik~$!A$ yang mungkin,
himpunan yang dihasilkan memuat $!A$ atau $\lnot !A$, dan (c)~setiap
kali himpunan itu memuat $!A \land !B$, himpunan tersebut juga memuat
$!A$ dan $!B$; jika memuat $!A \lor !B$, himpunan tersebut juga memuat
sekurang-kurangnya salah satu di antara $!A$ dan $!B$; dan seterusnya.
Kita terus melakukan hal ini (mungkin tanpa akhir). Sebut himpunan semua
!!{formula}s yang ditambahkan dengan cara tersebut~$\Gamma^*$. Konstruksi
di atas kemudian memberikan \iftag{FOL}{!!a{structure}~$\Struct{M}$}
{!!a{valuation}~$\pAssign{v}$} yang, melalui induksi, dapat kita buktikan
memenuhi semua kalimat dalam~$\Gamma^*$, dan karena itu juga semua
kalimat dalam~$\Gamma$, sebab $\Gamma \subseteq \Gamma^*$. Ternyata,
menjamin (a) dan~(b) sudah memadai. Himpunan kalimat yang memenuhi (b)
disebut \emph{lengkap}. Jadi, tugas kita ialah memperluas himpunan
konsisten~$\Gamma$ menjadi himpunan konsisten dan lengkap~$\Gamma^*$.

\iftag{FOL}{%
Rencana ini menghadapi satu kendala: jika $\lexists[x][!A(x)] \in
\Gamma$, kita berharap dapat memilih suatu !!{constant}~$c$ dan
menambahkan $!A(c)$ dalam proses tersebut. Namun, bagaimana kita tahu
bahwa hal itu selalu dapat dilakukan? Barangkali bahasa kita hanya
memiliki beberapa !!{constant}s, dan untuk setiap konstanta tersebut kita
memiliki $\lnot !A(c) \in \Gamma$. Kita tidak dapat sekaligus
menambahkan $!A(c)$ karena penambahan itu akan membuat himpunan tersebut
tak konsisten, dan kita tidak akan tahu apakah $\Struct{M}$ harus membuat
$!A(c)$ atau $\lnot !A(c)$ benar. Terlebih lagi, mungkin saja $\Gamma$
hanya memuat kalimat-kalimat dalam suatu bahasa yang sama sekali tidak
memiliki simbol konstanta (misalnya, bahasa teori himpunan).

Penyelesaian masalah ini adalah menambahkan tak hingga banyak konstanta
sejak awal, beserta kalimat-kalimat yang menghubungkan konstanta tersebut
dengan kuantor dengan cara yang tepat. (Tentu saja, kita harus memastikan
bahwa tindakan ini tidak dapat menimbulkan inkonsistensi.)

Konstruksi awal kita berfungsi dengan baik jika kalimat-kalimat atomik
hanya memuat !!{constant}s. Namun, bahasa tersebut mungkin juga memuat
!!{function}s. Dalam hal itu, menemukan fungsi-fungsi yang tepat
pada~$\Nat$ untuk ditetapkan kepada !!{function}s tersebut agar seluruh
konstruksi berhasil mungkin sulit. Karena itu, kita gunakan siasat lain: alih-alih
menggunakan $i$ untuk menafsirkan $\Obj c_i$, gunakan saja himpunan
!!{constant}s itu sendiri sebagai domain. Dengan demikian, $\Struct M$
dapat menetapkan setiap !!{constant} kepada dirinya sendiri:
$\Assign{\Obj c_i}{M} = \Obj c_i$. Namun, mengapa tidak melanjutkannya
sepenuhnya: tetapkan $\Domain{M}$ sebagai himpunan semua \emph{suku tertutup}
dalam bahasa tersebut!{} Dengan pilihan ini, terdapat penetapan fungsi
yang jelas bagi !!{function}s (fungsi-fungsi ini menerima suku tertutup sebagai
argumen dan menghasilkan suku tertutup sebagai nilai): kepada
!!{function}~$\Obj f^n_i$ kita tetapkan fungsi yang, ketika menerima
$n$ suku tertutup $t_1$, \dots,~$t_n$ sebagai masukan, menghasilkan suku tertutup
$\Obj f^n_i(t_1, \dots, t_n)$ sebagai nilainya.

Potongan teka-teki terakhir adalah cara menangani~$\eq$.
!!{predicate}~$\eq$ memiliki interpretasi tetap:
$\Sat{M}{\eq[t][t']}$ jika dan hanya jika $\Value{t}{M} =
\Value{t'}{M}$. Jika kita mengatur agar !!{value} suku~$t$ adalah $t$
sendiri, !!{structure} ini tidak akan membuat \emph{satu pun} kalimat
berbentuk $\eq[t][t']$ benar, kecuali jika $t$ dan $t'$ merupakan suku
yang sama persis. Tentu saja hal ini bermasalah, sebab hampir setiap
teori yang menarik dalam bahasa dengan !!{function}s memiliki kalimat
$\eq[t][t']$ sebagai teorema meskipun $t$ dan $t'$ bukan suku yang sama
(misalnya, dalam teori aritmetika: $\eq[(\Obj 0+ \Obj 0)][\Obj 0]$).
Untuk menyelesaikan masalah ini, kita ubah domain~$\Struct M$: alih-alih
menggunakan suku sebagai objek dalam~$\Domain{M}$, kita menggunakan
himpunan-himpunan suku, dan setiap himpunan memuat semua suku yang harus
disamakan menurut kalimat-kalimat dalam~$\Gamma$. Jadi, misalnya, jika
$\Gamma$ merupakan teori aritmetika, salah satu himpunan ini akan memuat
$\Obj 0$, $(\Obj 0 + \Obj 0)$, $(\Obj 0 \times \Obj 0)$, dan seterusnya.
Himpunan inilah yang kita tetapkan kepada $\Obj 0$, dan ternyata himpunan
ini juga merupakan nilai semua suku di dalamnya, misalnya nilai
$(\Obj 0 + \Obj 0)$. Oleh karena itu, kalimat
$\eq[(\Obj 0+ \Obj 0)][\Obj 0]$ akan benar dalam !!{structure} yang
telah diperbaiki ini.}{}

Jadi, inilah yang akan kita lakukan. Pertama, kita menyelidiki sifat-sifat
himpunan konsisten yang !!{complete}; khususnya, kita membuktikan bahwa
himpunan konsisten yang !!a{complete} memuat $!A \land !B$ jika dan hanya jika
himpunan itu memuat $!A$ dan~$!B$, memuat $!A \lor !B$ jika dan hanya
jika himpunan itu memuat sekurang-kurangnya salah satunya, dan seterusnya
(\olref[ccs]{prop:ccs}).\iftag{FOL}{ Kemudian kita mendefinisikan dan
  menyelidiki himpunan kalimat ``jenuh''. Himpunan jenuh adalah himpunan
  yang memuat implikasi-implikasi yang menghubungkan setiap !!{sentence}
  berkuantor dengan instans-instansnya
  (\olref[hen]{defn:henkin-exp}). Kita menunjukkan bahwa setiap himpunan
  konsisten~$\Gamma$ selalu dapat diperluas menjadi himpunan
  jenuh~$\Gamma'$ (\olref[hen]{lem:henkin}). Jika suatu himpunan
  konsisten, jenuh, dan !!{complete}, himpunan itu juga memiliki sifat
  bahwa ia memuat \iftag{prvEx}{$\lexists[x][!A(x)]$ jika dan hanya jika
    ia memuat $!A(t)$ untuk suatu suku tertutup~$t$}{}
  \iftag{defEx,defAll}{}{ dan }
  \iftag{prvAll}{$\lforall[x][!A(x)]$ jika dan hanya jika ia memuat
    $!A(t)$ untuk semua suku tertutup~$t$}{}
  (\olref[hen]{prop:saturated-instances}).}{}
Selanjutnya, kita mengambil himpunan konsisten yang\iftag{FOL}{ jenuh}{}
~$\iftag{FOL}{\Gamma'}{\Gamma}$ dan menunjukkan bahwa himpunan tersebut
dapat diperluas menjadi suatu \iftag{FOL}{himpunan yang jenuh, konsisten,
dan}{himpunan yang konsisten dan} !!{complete}~$\Gamma^*$
(\olref[lin]{lem:lindenbaum}). Himpunan $\Gamma^*$ inilah yang akan kita
gunakan untuk mendefinisikan \iftag{FOL}{model suku~$\Struct
  M(\Gamma^*)$}{!!{valuation}~$\pAssign v(\Gamma^*)$} kita.
\iftag{FOL}{Domain model suku tersebut adalah himpunan suku tertutup,
  sedangkan interpretasi !!{predicate}snya ditentukan oleh
  !!{sentence}s atomik}{Valuasi tersebut ditentukan oleh
  !!{propositional
    variable}s} dalam~$\Gamma^*$
(\olref[mod]{defn:termmodel}). Kita akan menggunakan sifat-sifat
himpunan konsisten yang\iftag{FOL}{ jenuh dan}{} lengkap untuk
menunjukkan bahwa
$\iftag{FOL}{\Sat{M(\Gamma^*)}}{\pSat{v(\Gamma^*)}}{!A}$ jika dan hanya
jika $!A \in \Gamma^*$ (\olref[mod]{lem:truth}), dan dengan demikian,
khususnya,
$\iftag{FOL}{\Sat{M(\Gamma^*)}}{\pSat{v(\Gamma^*)}}{\Gamma}$.
\iftag{FOL}{Terakhir, kita akan membahas cara mendefinisikan model suku
  jika $\Gamma$ juga memuat~$\eq$
  (\olref[ide]{defn:term-model-factor}) dan menunjukkan bahwa model itu
  memenuhi~$\Gamma^*$ (\olref[ide]{lem:truth}).}{}

\end{document}
